\documentclass{cv}

\usepackage[ngerman]{babel}
\usepackage{microtype}

% --- ATS/Upload safe: remove all clickable links/annotations ---
\makeatletter
\@ifpackageloaded{hyperref}{
  % hyperref is already loaded by the class -> disable link creation
  \let\href\@secondoftwo
  \let\url\@firstofone
  \hypersetup{draft}
}{}
\makeatother

\name{Marc Borer}
\address{Hirzengarten 2, CH-4226 Breitenbach}
\phone{+41~79~960~9691}
\email{mb@marc-borer.ch}
%\homepage{https://marc-borer.ch}
\linkedin{marc-b-6780ab5a/}
\github{karlhabsburg}

\begin{document}


\maketitle%

\section{Profil}
\setlength{\columnsep}{-5cm}
\begin{multicols}{2}
Senior ICT Systems Engineer – Spezialist für Microsoft Cloud-Infrastrukturen (Azure, M365) und DevOps-Automatisierung. Eigenverantwortliche Durchführung komplexer Migrationsprojekte (Exchange Online, Teams Calling, Hybrid Identity) sowie Entwicklung automatisierter Compliance- und Lifecycle-Management-Plattformen.\\ Über 10 Jahre Erfahrung in der Konzeption und Umsetzung unternehmenskritischer IT-Systeme mit Fokus auf Skalierbarkeit, Security und Prozessoptimierung.\\
    \includegraphics[width=3.4cm, right]{profile.png}
\end{multicols}

\section{Berufs-\\erfahrung}
\subsection{CTC Analytics AG}[Zwingen, CH]
\begin{positions}
    \entry{IT Systems Engineer}{Mar 2019~--~Heute}
\end{positions}
Verantwortlich für Architektur, Weiterentwicklung und sicheren Betrieb der Microsoft Cloud- und Hybrid-Infrastruktur eines international tätigen Unternehmens im Bereich Laborautomation mit ca. 180 Mitarbeitenden. Betreuung von über 100 Windows- und Linux-Servern mit Schwerpunkt auf Microsoft 365, Exchange Online, Teams, Active Directory sowie Infrastruktur-Automatisierung und Standardisierung.

\textbf{Schlüsselprojekte und Verantwortlichkeiten:}

\begin{itemize}[leftmargin=*]

\item \textbf{Cloud-Migration Exchange Online \& Microsoft Teams}

Projektlead und technische Gesamtverantwortung für die Migration der Unternehmenskommunikation (300+ Mailboxen, globale Telefonie) von Exchange 2013 und Skype for Business zu Microsoft 365.

Durchführung von Ist-/Soll-Analyse, Zielarchitektur-Definition sowie Erstellung einer Migrations-Roadmap inkl. Best-/Middle-/Worst-Case-Szenarien und Timeline-Steuerung. Stakeholder-Management über Fachbereiche, Management und externe Partner hinweg.

\smallskip
\textbf{Projektumfang:}
\begin{itemize}[nosep]
    \item Architekturdesign Hybrid-Koexistenz Exchange Online/Exchange Server 2019
    \item Modernisierung der Voice-Infrastruktur inkl. SBC und SIP-Trunk-Integration
    \item Microsoft Teams Calling Rollout inkl. Direct Routing
    \item Migration von 300+ Mailboxen ohne Serviceunterbruch
\end{itemize}



\vspace{0.3cm}

\item \textbf{Enterprise Device Inventory \& Compliance Platform}

Konzeption und Umsetzung einer vollautomatisierten BIOS- und Compliance-Lifecycle-Plattform für 500+ Intune-verwaltete Endgeräte auf Basis einer serverlosen Azure-Architektur.

\smallskip
\textbf{Komponenten:}
\begin{itemize}[nosep]
    \item \textbf{Monitoring:} Azure Functions und Proactive Remediation Scripts zur täglichen Erfassung von BIOS-Versionen, Applikations-Inventory und BitLocker-Status via Log Analytics
    \newpage
    \item \textbf{Deployment:} Azure Automation Runbooks mit dynamischer OEM-Katalog-Integration (Lenovo/Dell/HP), automatisierter Intune Win32-Paket-Generierung sowie BitLocker Suspend/Resume-Orchestrierung
    
    \item \textbf{Orchestration:} Dynamische Azure AD-Gruppen, ring-basiertes Rollout \\(Canary→Pilot→Broad) sowie Package Supersedence via Graph API
    
    \item \textbf{Analytics:} Azure Monitor Workbooks und Logic Apps zur Erstellung von Compliance-Dashboards und proaktiven Alerting-Workflows
\end{itemize}
%\newpage
\vspace{0.3cm}

\item \textbf{Hybrid Identity \& User Lifecycle Automation}

Design und Implementierung einer hybriden Automatisierungslösung für vollständig standardisiertes Mitarbeiter-Onboarding und -Offboarding.

\smallskip
\textbf{Technische Umsetzung:}
\begin{itemize}[nosep]
    \item HR-initiiertes Self-Service-Portal mittels Power Automate und PowerApps
    \item PowerShell-basierte Orchestrierung für Active Directory Account-Provisionierung,\\Exchange Online Mailbox-Erstellung, Microsoft Teams-Lizenzierung, Fileshare-\\Berechtigungen sowie Microsoft 365-Gruppenmitgliedschaften
\end{itemize}

\vspace{0.3cm}

\item \textbf{SIEM-Implementation mit Elastic Stack}

Aufbau einer unternehmensweiten Security Information and Event Management-Plattform zur zentralisierten Protokollierung und Security-Analyse.

\smallskip
\textbf{Implementierung:}
\begin{itemize}[nosep]
    \item Elasticsearch-Logstash-Kibana Cluster (High Availability Setup)
    \item Strukturierte Multi-Source Data Ingestion via Filebeat (Application Logs), Metricbeat (Performance Metrics), Winlogbeat (Windows Events), Heartbeat (Uptime Monitoring)
    \item Entwicklung von Custom Kibana Dashboards für Security Operations (Threat Detection, Anomaly Analysis, Compliance Reporting)
\end{itemize}


\end{itemize}

\subsection{Sabbatical}[]
\begin{positions}
    \entry{Persönliche Auszeit}{Jan 2018 - Feb 2019}
\end{positions}
%\newpage
\subsection{Vision Information Transaction AG / Picturepark}[Aarau, CH]
\begin{positions}
    \entry{IT Operations Manager}{Mai 2014~--~Dez 2017}
\end{positions}
Mitarbeiter im IT-Operations-Team eines international tätigen SaaS-Anbieters. Verantwortlich für die Entwicklung und den Betrieb der Picturepark-Plattform sowie kundenspezifische Implementierungen.

\textbf{Hauptverantwortlichkeiten:}

\begin{itemize}
    \item \textbf{CI/CD-Pipeline für Azure-Deployments}\\Konzeption und Implementierung automatisierter Deployment-Prozesse für Alpha- und Beta-Versionen der Picturepark-Software auf Microsoft Azure. Technologie-Stack umfasste Ansible, Bamboo, sowie automatisierte Installation von Active Directory, DNS, IIS, HAProxy und ELK-Cluster.
    %\item \textbf{SaaS-Infrastruktur Design und Betrieb}\\Aufbau und Administration einer hochverfügbaren SaaS-Infrastruktur basierend auf Hyper-V-Clustern mit SAN-Symphony-Storage, SQL Server High-Availability-Konfiguration und HAProxy-Load-Balancing.
    \item \textbf{Kundenimplementierungen}\\Technische Leitung bei der Installation und Konfiguration der Picturepark-Software bei Enterprise-Kunden auf Multi-Server-Umgebungen.
\end{itemize}
\newpage
\section{Ausbildung}
\subsection{Wirtschaftsinformatikschule}[Zürich, CH]
\begin{positions}
    \entry{Informatiker EFZ Schwerpunkt Systemtechnik}{2010~--~2014}
\end{positions}

\section{Kurse}
\subsection{Digicomp}[Zürich, CH]
\begin{positions}
    \entry{Designing and Implementing MS DevOps Solutions}{4 Tage, 2025}
    \entry{PowerShell – Fortgeschrittene Techniken}{2 Tage, 2019}
\end{positions}

\section{Technische\\Kompetenzen}
\begin{description}
\item[Cloud \& Identity] Microsoft Azure, Hybrid Identity, Azure AD / Entra ID, ADFS, WAP
\item[Microsoft 365 \& Collaboration] Exchange Online \& On-Premises, Microsoft 365, Microsoft Teams, Teams Calling (Direct Routing)
\item[Endpoint \& Lifecycle Management] Microsoft Intune (MDM/MAM), Client- \& Server Patch Management, Application Deployment (Altiris)
\item[Automation \& Scripting] PowerShell (Experte), CI/CD-gestützte Automatisierung
\item[Virtualisierung \& Plattformen] Hyper-V, VMware vSphere / ESXi
\item[Monitoring \& Logging] ELK Stack (Elasticsearch, Logstash, Kibana), SCOM, Nagios
\item[Netzwerk \& Security] Load Balancing, HAProxy, VPN, Security-nahe Betriebsprozesse
\item[Betriebssysteme] Windows Server (2008–2022), Red Hat Enterprise Linux, CentOS
\item[Weitere Technologien] SQL Server, IIS, Bash-Scripting
\end{description}

%\section{Zertifikate}
%\textbf{Microsoft Certified: DevOps Engineer Expert} \printdate{2025} \\ 

\section{Persönliche \\Angaben}
\textbf{Geburtsdatum} 22. Januar 1992\\
\textbf{Familienstand} Ledig\\
\textbf{Sprachen\\} 
\hspace*{0.4cm}Deutsch (Muttersprache)\\
\hspace*{0.4cm}Englisch (C2 - Sehr gute Kenntnisse)\\
\textbf{Nationalität} Schweizer \\
\textbf{Interessen} Turnverein, Klavier, Wandern, Skifahren \\
\textbf{Verfügbarkeit} 3 Monate Kündigungsfrist

\section{Referenzen}
\begin{description}
    \item[Auf Anfrage verfügbar]
\end{description}

\end{document}
